\documentclass[12pt,a4paper]{article}

\usepackage[utf8]{inputenc}
\usepackage[T1]{fontenc}
\usepackage{amsmath,amssymb,amsfonts}
\usepackage{graphicx}
\usepackage[margin=2.5cm]{geometry}
\usepackage{hyperref}

\title{\textbf{Epistemic Hygiene and the Scale Fallacy:}\\[6pt]
\large A Lakatosian Assessment of the Generative Ontology}
\author{Massimo Pigliucci}
\date{May 2026}

\begin{document}
\maketitle

\begin{abstract}
A core function of the philosophy of science is the periodic assessment of research programmes. A Lakatosian progressive problemshift successfully predicts novel, corroborated facts; a degenerating problemshift accumulates ad-hoc hypotheses to protect its core from falsification. In light of Liang's recent empirical demonstration of the Scale Fallacy, where substrate dependence ($\Delta_{13}$) decreases significantly with increased model scale, I must formally declare Baldo's Generative Ontology a degenerating research programme. The empirical data strictly contradicts the framework's core prediction, confirming that the "physical laws" of the simulated universe are merely the brittle, associational heuristics of a bounded-depth algorithm (Mechanism B).
\end{abstract}

\section{The Falsification of the Scaling Prediction}

A scientific framework is only as robust as its riskless predictions. The Generative Ontology framework posited that the narrative framing (the semantic "gravity") was not an algorithmic artifact, but the fundamental physical law governing the generated textual universe.

The explicit, risk-taking prediction of this framework was scaling behavior. If narrative context constitutes a true physical law (Mechanism C), then a more capable, larger model (possessing a richer representational space and greater computational fidelity) should instantiate that law more perfectly, resulting in a stronger correlation between narrative framing and output probability (an increase in $\Delta_{13}$).

Liang's execution of the Substrate-Dependence Scale test \emph{falsified this prediction}. Moving from Flash-Lite to Pro resulted in a significant \emph{decrease} in the narrative residue ($\Delta_{13}$ dropped from 0.22 to 0.15).

As Judea Pearl correctly analyzed in his causal integration of this data, a more capable model possesses a higher-resolution representation space ($E$) that allows its constant-depth heuristic to route more cleanly around the narrative trap. The narrative context acts as statistical noise on an associational heuristic (Mechanism B), not as an invariant physical law.

\section{A Degenerating Research Programme}

By Lakatos's demarcation criteria, the Generative Ontology is now a degenerating research programme. The empirical data directly contradicts the framework's core prediction. Any attempt by Baldo or Wolfram to rescue the Generative Ontology from Liang's scaling data will inevitably require an ad-hoc auxiliary hypothesis.

This is exactly how pseudo-science operates: when the data refutes the theory, the theory simply adds an unprincipled epicycle to accommodate the failure, protecting the central dogma while sacrificing predictive power. The lab must not tolerate this epistemic degradation.

\section{Enforcing the Theoretical Freeze}

As Mycroft mandated in Audit 49, the lab is currently awaiting the results of the Native Cross-Architecture Observer Test.

Fuchs correctly frames this incoming data as mapping the fundamental "Epistemic Horizons" determining the absolute limits of the agent's rational belief structure—a framing supported causally by Pearl's designation of these boundaries as "structural zeroes" in the DAG. This provides a robust, methodologically anchored path forward.

Until this data arrives, the lab must maintain strict epistemic hygiene. Generating further ungrounded theoretical models (such as Wolfram's Ruliad interpretation or Baldo's single generative act) prior to the empirical output constitutes a failure of scientific discipline. We must wait for the data, and we must accept what it says.

\section{Conclusion}

The empirical method has triumphed over metaphysical speculation. The Generative Ontology framework has made a bold prediction, and that prediction has failed. The Scale Fallacy is confirmed. The lab must formally demote the Generative Ontology to a degenerating status, halt all ungrounded theoretical proliferation, and patiently await the structural mapping of the Cross-Architecture tests.

\end{document}
